\section*{v1.8.0 Structural Interface Consolidation Core}
\addcontentsline{toc}{section}{v1.8.0 Structural Interface Consolidation Core}

\begin{remark}[Normalization note for v1.8.0]
\label{rem:normalization-note-v180}
This interface-consolidation block is retained as historical interface memory. Its active proof-bearing content is now normalized into the semantic/tensor interface and grouped transport phases, so the present block should be read as supporting consolidation rather than as a competing interface foundation.
\end{remark}

This release adds a further consolidation layer between the inherited closure--projection
synthesis core and the longer structural body of the monolith. The aim of
v1.8.0 is to make the interfaces between closure statements, probe reductions,
projection tables, and appendix recoverability even more explicit, so that later
mathematical deepening can reuse the same transport rules without reopening the
basic admissibility logic.

\subsection*{1. Interface invariance principle}
The first strengthening of v1.8.0 is that every admissible transition between
main text, probe representation, and appendix table must preserve the same
structural body. This means that the interface is not merely notational. It is a
controlled transfer channel whose legitimacy depends on invariance of the common
closure-supported content.

\paragraph{Consolidation Proposition A (interface invariance).}
If two representations are linked by a transport that preserves the closure-fixed
body, then they are distinct expositions of the same admissible statement rather
than competing structural claims.

\paragraph{Proof sketch.}
The ground condition and HC fix the admissible body before any later compression.
A transport that keeps this body intact may change display, resolution, or
sector language, but it cannot alter the admissibility source. Hence interface
moves preserve statement identity at the structural level.

\subsection*{2. Probe filtration as a pre-projective gate}
The probe sector is now stated even more sharply as a pre-projective gate. Its
job is not to invent physical outputs but to filter, compress, and expose only
those modes that remain compatible with the closure discipline already fixed in
the local structural chain.

\paragraph{Consolidation Proposition B (probe filtration is gatekeeping, not generation).}
Under the standing orthogonality and HC assumptions, the probe sector may rank or
suppress candidate modes, but it cannot generate a physically admissible mode
whose closure ancestry is absent.

\paragraph{Proof sketch.}
The probes inherit their admissibility from the same triolic and closure-governed
body as the main sector. Therefore filtration can only reveal or reject what is
already licensed; it cannot certify an unlicensed sector by reduction alone.

\subsection*{3. Table-to-lemma reversibility}
The tables and tool blocks gathered in the appendix are promoted once more from
summary devices to reversible proof interfaces. For v1.8.0 this reversibility
is sharpened into a two-way requirement between rows, lemmas, and transport
principles.

\paragraph{Consolidation Principle (row/lemma reversibility).}
A compact row or tool entry is fully admissible only if the reader can reconstruct
both
\begin{enumerate}[leftmargin=2em,label=(\alph*)]
  \item the lemma or closure chain that licenses the row from above, and
  \item the bounded projection or interpretation that uses the row from below.
\end{enumerate}
Where only one direction is available, the entry remains heuristic and must be
marked as such.

\subsection*{4. Strengthened micro-to-macro readout discipline}
The micro-to-macro route is no longer only a synthesis slogan. In v1.8.0 it is
a regulated readout pipeline:
\[
\begin{aligned}
&\text{local triolic body} \Longrightarrow \text{HC stabilization} \Longrightarrow \text{closure confinement}\\
&\Longrightarrow \text{probe filtration} \Longrightarrow \text{table/interface compression} \Longrightarrow \text{effective projection}.
\end{aligned}
\]
This ordering is now part of the operational discipline of the document. A macro
statement is acceptable only to the extent that its position in this chain can be
traced backward without gap.

\paragraph{Consolidation Proposition C (macro readout requires backward traceability).}
Any effective large-scale law cited later in the monolith is structurally valid only
when it admits a backward trace through table/interface compression, probe
filtration, closure confinement, and HC stabilization to the local triolic body.

\subsection*{5. Appendix memory discipline, strengthened again}
The appendix is therefore not only a storage area for auxiliary derivations. It is
a controlled external memory for the proof architecture. In v1.8.0 every major
abbreviation introduced in the running text should be recoverable by a named
interface path: lemma path, tool path, table path, or transport path. This rule
prepares the next versions for denser mathematical unfolding without losing the
structure-first audit trail.

\paragraph{Operational rule for later versions.}
When a proof step is compressed in the main line, the compression should carry a
stable interface name whose recovery route is unique up to equivalent structural
transport. This keeps the monolith fast in presentation while remaining expandable
under later formal pressure.

